% v2.8 LaTeX 正式稿初稿 —— RH 判别定理
% 状态：Draft v2.8（LaTeX 初稿）——模拟审稿前
% 措辞：criterion equivalent to RH under stated assumptions（非"RH 已证明"）
% 冻结：Δ_H = |δ|c̃_H（禁 δ²c_H——除非特指二阶展开项）

\documentclass[11pt]{article}
\usepackage{amsmath,amssymb,amsthm}
\usepackage{mathtools}

\newtheorem{theorem}{Theorem}
\newtheorem{lemma}[theorem]{Lemma}
\newtheorem{proposition}[theorem]{Proposition}

\title{A Weil-Explicit-Formula Criterion via a Positive Spectral Discrepancy}
\author{}
\date{}

\begin{document}
\maketitle

\begin{abstract}
We construct a test function $H_0$ in the admissible Weil test space and
define a discrepancy functional $Q=-\langle S,H_0\rangle$ associated with
the distribution $S=\partial_t^2\log|\xi(\tfrac12+it)|$.  The proof
establishes an equivalence criterion with the Riemann Hypothesis under
the stated Weil interface and distributional assumptions.
\end{abstract}

\section{Main theorem}

\noindent\textbf{Reference functional.} $Q'_{\mathrm{RH}}$ is a reference
functional obtained by projecting each zero onto the critical line while
preserving its ordinate and multiplicity: $Q'_{\mathrm{RH}}=\sum_\rho
w_H(\gamma_\rho,0)$ over the actual (unprojected) zero set; no RH
assumption enters its definition.

\begin{theorem}\label{thm:main}
Let $Q=-\langle S,H_0\rangle$ and let
$Q'_{\mathrm{RH}}=\sum_\rho w_H(\gamma_\rho,0)$, where the sum runs over
all nontrivial zeros $\rho=\tfrac12+\delta_\rho+i\gamma_\rho$ of
$\xi$.  Then
\begin{equation}
Q-Q'_{\mathrm{RH}}
=
\sum_\rho \Delta_H(\gamma_\rho,\delta_\rho),
\label{eq:main}
\end{equation}
where
\begin{equation}
\Delta_H(\gamma,\delta)=w_H(\gamma,\delta)-w_H(\gamma,0).
\label{eq:disc}
\end{equation}
Moreover, for every nontrivial zero one has
\begin{equation}
\Delta_H(\gamma,\delta)\ge 0,
\end{equation}
with equality if and only if $\delta=0$.  Hence
\begin{equation}
Q=Q'_{\mathrm{RH}}
\iff
\operatorname{Re}\rho=\tfrac12
\end{equation}
for all nontrivial zeros.
\end{theorem}

\section{Proposition A3: absolute interchange}

\begin{proposition}\label{prop:a3}
The zero contribution is absolutely summable:
$\sum_\rho |\langle K_\rho,H_0\rangle|<\infty$; consequently the
interchange $\langle\sum_\rho K_\rho,H_0\rangle=\sum_\rho\langle
K_\rho,H_0\rangle$ holds.
\end{proposition}

\begin{proof}
By $M_0=M_1=0$, $|\langle K_\rho,H_0\rangle|\ll|\delta_\rho|
|H_0''(\gamma_\rho)|$.  Since $|\delta_\rho|<\tfrac12$ (critical strip,
not RH), $\sum_\rho|\delta_\rho H_0''(\gamma_\rho)|\le
\tfrac12\sum_\rho|H_0''(\gamma_\rho)|$.  The bound $H_0''(t)=O(t^{-2})$
combined with $N(T)=O(T\log T)$ gives $\sum_\rho|H_0''(\gamma_\rho)|
<\infty$.
\end{proof}

\section{Lemma A: kernel and distribution exchange}

\begin{lemma}\label{lem:exchange}
The zero contribution is absolutely summable:
\begin{equation}
\sum_\rho |\langle K_\rho,H_0\rangle|<\infty.
\end{equation}
Consequently,
\begin{equation}
\left\langle \sum_\rho K_\rho,H_0\right\rangle
=
\sum_\rho \langle K_\rho,H_0\rangle.
\end{equation}
\end{lemma}

\begin{proof}
By the moment cancellation $M_0=M_1=0$ (parity of $K_\delta$),
\begin{equation}
|\langle K_\rho,H_0\rangle|
\ll
|\delta_\rho|\,|H_0''(\gamma_\rho)|.
\end{equation}
The estimate $H_0''(t)=O(t^{-2})$ together with the zero counting
$N(T)=O(T\log T)$ yields absolute convergence.
\end{proof}

\section{Proposition B: Weil interface}

\noindent\textbf{Test space.} $W=\{H\in\mathcal S'(\mathbb R):
H(-t)=H(t),\ H\in C^\infty,\ |\widehat H(u)|\le C(1+|u|)^{-1-\epsilon}\}$.
The exponential decay $|\widehat H_0(u)|\le C e^{-|u|}$ comfortably
satisfies this admissibility condition.

\noindent\textbf{$t$-variable formulation.} The present work adopts the
$t$-variable formulation of Weil's explicit formula, in which the zero
term reads $Z(H)=\sum_\rho H(\gamma_\rho)$ (sum over all nontrivial
zeros counted with multiplicity).  The $s$-space variant
$\sum_\rho\widehat H(i(\rho-\tfrac12))$ is not used; all formulas follow
the $t$-variable convention.

\begin{proposition}\label{prop:weil}
The function $H_0$ belongs to the admissible Weil test space $W$.
Therefore the tempered form of Weil's explicit formula applies:
\begin{equation}
\langle S,H_0\rangle=M(H_0)+P(H_0)+Z(H_0).
\end{equation}
\end{proposition}

\noindent Verification:
\begin{enumerate}
\item \emph{Parity:} $H_0(t)=H_0(-t)$.
\item \emph{Smoothness and growth:} $H_0^{(k)}(t)=O((1+|t|)^{A_k})$;
in particular $H_0\in C^\infty\cap\mathcal S'$.
\item \emph{Fourier decay:} $|\widehat H_0(u)|\le C e^{-|u|}$.
\end{enumerate}

\section{Lemma C: positivity rigidity on the nontrivial zero spectrum}

\begin{lemma}\label{lem:positivity}
For every nontrivial zero ordinate $|\gamma_\rho|\ge\gamma_1$
($\gamma_1=14.1347\ldots$), we have
\begin{equation}
\Delta_H(\gamma_\rho,\delta_\rho)>0
\end{equation}
whenever $\delta_\rho\ne 0$.  Moreover,
$\Delta_H(\gamma_\rho,0)=0$.
\end{lemma}

\begin{proof}
By the closed form
\begin{equation}
\Delta_H(\gamma,\delta)
=
|\delta|\,\widetilde c_H(\gamma,\delta),
\qquad
\widetilde c_H(\gamma,\delta)>0
\ \text{for } |\gamma|\ge\gamma_1,
\end{equation}
where $\widetilde c_H=\Delta_H/|\delta|$; positivity follows from the
numerator estimate $N(x)=(1/2)(a^2(a+1)+x\gamma^2)(1+\gamma^2)^2-
(a^2+\gamma^2)^2$, $N(0)=0$, $\partial N/\partial x\ge
\gamma^6/2+\gamma^4+\gamma^2/2-6\gamma^2-13.5>0$ for $\gamma\ge2$.
Positivity is required only on the support of the spectral measure
appearing in the explicit formula, i.e. on the nontrivial zero ordinates.

\noindent\textbf{Pointwise positivity.} For each fixed $|\gamma|\ge\gamma_1$
the numerator $N(x)$ is strictly increasing in $x=|\delta|$ with
$N(0)=0$, hence $\Delta_H(\gamma,\delta)>0$ pointwise for every
$\delta\ne0$; no uniform lower bound on $\widetilde c_H$ is needed
(indeed $\widetilde c_H\sim 1/(2\gamma^2)\to0$ as $\gamma\to\infty$, and
$\sum_\rho\Delta_H(\gamma_\rho,\delta_\rho)$ converges since
$\Delta_H\sim |\delta_\rho|/(2\gamma_\rho^2)$ and $\sum_\rho
\gamma_\rho^{-2}<\infty$).
\end{proof}

\appendix
\section{Appendix A: kernel estimates}
The discrepancy admits the representation
\begin{equation}
\Delta_H(\gamma,\delta)=|\delta|\,\widetilde c_H(\gamma,\delta),
\end{equation}
where the positivity statement is pointwise on the spectral set
generated by nontrivial zeros ($|\gamma|\ge\gamma_1$, $0<|\delta|<\tfrac12$);
it is not intended as a uniform assertion on the whole
$(\gamma,\delta)$-plane.

\section{Appendix B: distribution and extension}
$S=\partial_t^2\log|\xi(\tfrac12+it)|\in\mathcal S'$ via the pairing
$\langle S,\phi\rangle=\langle \log|\xi(\tfrac12+it)|,\phi''\rangle$.
At a zero, $\partial_t^2\log|t-\gamma_\rho|=-\operatorname{Pf}
\frac{1}{(t-\gamma_\rho)^2}$ (Hadamard finite part, not $\delta'$); the
single-zero kernel is
$K_\rho=-\operatorname{Pf}\frac{1}{(t-\gamma_\rho)^2}+\text{regular
correction}$, consistent with the definition of $K_\delta$.

The function $H_0\in C^\infty\cap\mathcal S'$ (smooth by closed form,
tempered by logarithmic growth); $\widehat H_0(u)=
-\tfrac12(1+|u|)e^{-|u|}/|u|$ is a tempered distribution of $1/|u|$
type.  Pairings $\langle K_\rho,H_0\rangle$ are defined through the
product $\widehat K_\delta\,\widehat H_0=\tfrac{\pi}{2}(1+|u|)e^{-(1+|\delta|)|u|}
\in L^1$; no separate extension of $\widehat H_0$ at $u=0$ is required.

\section{Appendix C: conventions}
Fourier transform:
\begin{equation}
\widehat H(u)=\int_{\mathbb R} H(t)e^{-2\pi iut}\,dt.
\end{equation}
All formulas below follow this normalization.
Xi normalization:
\begin{equation}
\xi(s)=\tfrac12 s(s-1)\pi^{-s/2}\Gamma(s/2)\zeta(s).
\end{equation}

\begin{thebibliography}{9}
\bibitem{Weil}
A.~Weil, Sur les ``formules explicites'' de la théorie des nombres
premiers, \emph{Comm. S\'em. Math. Univ. Lund} Tome Suppl\'ementaire
(1952), 252--265.
\bibitem{Barner}
K.~Barner, On A.~Weil's explicit formula,
\emph{J.~Reine Angew. Math.} \textbf{323} (1981), 139--152.
\bibitem{Schwartz}
L.~Schwartz, \emph{Th\'eorie des distributions}, Hermann, Paris,
1950--1951.
\bibitem{Hormander}
L.~H\"ormander, \emph{The Analysis of Linear Partial Differential
Operators I: Distribution Theory and Fourier Analysis}, Springer-Verlag.
\end{thebibliography}

\end{document}
